% Includable research note. Requires amsmath, amssymb, amsthm.
% Uses theorem, lemma, corollary, proposition, proof environments.
% This is a standard-N result; it is not a positive iPRA axiomatization.

\section{Strict dependence in the positive-iteration term algebra}
\label{sec:positive-strict-dependence}

\paragraph{Status and scope.}
The current checker-and-barriers draft proves a zero-profile obstruction
for positive iteration. The certified-zero draft gives the stronger,
whole-theory obstruction of a two-element model for the proposed eager
arithmetic, even with induction on every formula. The results below concern
a different question: the exact functions denoted by the positive-iteration
terms in the \emph{standard natural numbers}. They strengthen the
standard-function analysis, but do not remove that finite-model obstruction
or establish a smaller iPRA presentation. A search of the supplied TeX
corpus found no statement of the strict-dependence classification below;
no literature-wide priority claim is made.

\paragraph{The term algebra.}
Let terms be built from variables, zero, successor, and
\[
 F(n,a;u.h)=h_{\vec x}^{\,n+1}(a),\qquad
 h_{\vec x}(b)=h(b,\vec x).
\]
Only $u$ is bound in the update $h$. All passive parameters are fixed
during iteration, the counter is not passed to the update, and each
description is finite earlier-stage syntax. Arbitrary nesting and
capture-free substitution are allowed. Equivalently,
\[
 F(0,a;u.h)=h[a/u],\qquad
 F(Sn,a;u.h)=h[F(n,a;u.h)/u].
\]
The following assertions are about this standard interpretation, not about
every model of those equations.

For a function $f:\mathbb N^r\to\mathbb N$, a coordinate $j$ is
\emph{essential} if changing that coordinate can change the value, with
the other coordinates fixed. Write $U(f)$ for its set of essential
coordinates. Strict increase in coordinate $j$ means strict increase for
\emph{every} fixed choice of the other coordinates, not merely for some
such choice.

\begin{theorem}[Constant--projection--growth classification]
\label{thm:positive-classification}
Every fixed positive-iteration term $t(x_1,\ldots,x_r)$ belongs to exactly
one of the following three classes.
\begin{enumerate}
\item It is a constant function.
\item It is one coordinate projection $x_j$.
\item Its essential-coordinate set $U$ is nonempty; it is independent of
      the remaining coordinates; it is strictly increasing in each
      coordinate in $U$; and
      \[
             t(\vec x)\geq x_j+1\qquad(j\in U).
      \]
\end{enumerate}
In particular every term is coordinatewise nondecreasing. Its class and
its exact essential-coordinate set can be calculated by structural
recursion on its finite syntax. Calculating the class does not require
evaluating a constant term or uniformly evaluating arbitrary primitive
recursive descriptions.
\end{theorem}

\begin{proof}
We prove the classification, monotonicity, and the asserted exact
dependencies together. Zero is constant and a variable is its projection.
Successor preserves constancy for a constant argument. For a nonconstant
argument it preserves its exact dependencies and strict increases, and
raises its lower bound from at least $x_j$ to at least $x_j+1$.

Consider $F(n,a;u.h)$. Apply the induction hypothesis to $n$, to $a$, and
to $h$, regarding its state $u$ as an additional independent coordinate.
There are four immediate cases.
\begin{enumerate}
\item If $h$ is constant, the result is that constant: at least one
      update is performed.
\item If $h$ is the state projection $u$, the result is the seed $a$.
\item If $h$ is a passive-coordinate projection, the result is that
      projection, independently of the seed and counter.
\item If $h$ has the third class but does not depend on $u$, the result
      is $h(\vec x)$, with exactly its already classified dependencies.
\end{enumerate}
The remaining case is that $h$ has the third class and depends on $u$.
For every fixed passive assignment, therefore,
\[
                 h(b,\vec x)\geq b+1,
\]
and $h$ is strictly increasing in $b$. Let $U_n,U_a$ be the essential
sets of the counter and seed, and let $V$ be the passive essential set of
$h$. Define
\[
                         U=U_n\cup U_a\cup V.
\]
The resulting function is independent of coordinates outside $U$.
If $U$ is empty it is constant. Suppose $U$ is nonempty.

First, the result is monotone in every coordinate. For two ordered
parameter assignments $\rho\leq\sigma$, the seed and counter values
are ordered, and the update under $\rho$ is pointwise at most the update
under $\sigma$. Induction over any common number of updates preserves
the ordered states. Any additional updates at $\sigma$ can only increase
its state, because each update raises the state by at least one.

Next, every coordinate in $U$ gives a strict increase. An increase of
an essential seed coordinate strictly increases the initial state, and
strict monotonicity of the update in its state preserves that difference.
An increase of an essential passive update coordinate strictly increases
the state at the first update, which exists even at counter zero; later
updates preserve the difference. An increase of an essential counter
coordinate increases the number of updates, and every extra update
raises the state by at least one. If a coordinate occurs in more than
one of these roles, all its other effects are nondecreasing, so the
same strict comparison remains valid. Thus $U$ is precisely the
essential-coordinate set.

Finally, writing $N,A$ for the counter and seed values, the result is at
least $A+N+1$. If $j\in U_a$, the induction hypothesis gives
$A\geq x_j$, so the result is at least $x_j+1$. If $j\in U_n$, it
gives $N\geq x_j$, with the same conclusion. If $j\in V$, the last
application of $h$ is at least $x_j+1$ by the induction hypothesis on
$h$. This proves the required dominance for every $j\in U$.

These cases also constitute a finite classification algorithm. Constants
and projections have their explicit tags. In the last case one forms
the union $U_n\cup U_a\cup V$ and tests whether it is empty. No
semantic search or evaluation of all assignments is involved.
The three outcomes are disjoint: a projection fails the strict additive
lower bound at that coordinate, and a constant has no essential
coordinate.
\end{proof}

\begin{corollary}[No bounded nonconstant numerical term]
Every nonconstant positive-iteration term is unbounded in each of its
essential coordinates with all other coordinates fixed. Consequently no
term denotes a nonconstant function with finite range. In particular
there is no Boolean-valued zero-test term, and no predecessor term.
\end{corollary}

\begin{proof}
A projection is unbounded in its sole essential coordinate. In the third
class the lower bound $t(\vec x)\geq x_j+1$ proves unboundedness.
Predecessor is nonconstant but is not strictly increasing from $0$ to $1$.
\end{proof}

\subsection{An obstruction for tuple domains with actual equality}

The classification also obstructs moving to finite tuples if interpreted
equality remains actual coordinatewise equality. Quotient equality is a
material exception, as the next subsection demonstrates explicitly.

\begin{theorem}[No nontrivial term-defined unital magma on a tuple domain]
\label{thm:positive-actual-unit}
Let $d$ be finite, let $\varnothing\ne D\subseteq\mathbb N^d$, and let
$e\in D$. Suppose $f:D\times D\to D$ is given by a fixed tuple of
positive-iteration terms, with any fixed numerical parameters, and
\[
                         f(e,z)=z=f(z,e)\qquad(z\in D)
\]
in actual coordinatewise equality. Then $D$ has exactly one element.
No definability assumption on $D$ is needed.
\end{theorem}

\begin{proof}
Fixed numerical parameters can be replaced by their numeral terms, so
Theorem~\ref{thm:positive-classification} applies to every coordinate
function $f_i(x_1,\ldots,x_d,y_1,\ldots,y_d)$.
Let $J$ be the coordinates that are nonconstant on $D$. If $i\in J$,
$f_i$ must depend on at least one coordinate from each input tuple:
otherwise one of $f_i(e,z)$ and $f_i(z,e)$ would be constant on $D$.
In particular $f_i$ is neither constant nor a projection, and has the
third class.

There is an essential left-input coordinate $j_i\in J$. Indeed, if
every essential left-input coordinate were constant on $D$, then
$f_i(z,e)$ would be constant there, contrary to $i\in J$.
The dominance conclusion and the right unit law yield
\[
                              z_i\geq z_{j_i}+1
                              \qquad(z\in D).
\]
If $J$ were nonempty, the map $i\mapsto j_i$ on the finite set $J$
would have a directed cycle. Summing the displayed inequalities around
that cycle gives $z_i\geq z_i+\ell$ for its positive length $\ell$,
an impossibility. Thus $J$ is empty. Every coordinate is constant on
$D$, so the nonempty set $D$ is a singleton. The case $d=0$ is immediate.
\end{proof}

\subsection{The same algebra does have an injective pairing term}

The preceding theorem is not an almost-unary majorization result.
Positive iteration combines two unbounded inputs, and even defines an
injective numerical pairing in the standard model.
Put
\[
 q(a,b):=F(a,b;u.Su)=a+b+1,\qquad
 d(a):=q(a,a)=2a+1,
\]
and define
\[
             \pi(a,b):=F(a,S(d(b));u.d(u)).
\]
Iteration of $u\mapsto2u+1$ gives the exact identity
\[
             \pi(a,b)+1=2^{a+1}(2b+3).
\]
The factor $2b+3$ is odd. Uniqueness of the power of two dividing a
positive integer therefore recovers $a$ from $\pi(a,b)$; the odd part
then recovers $b$. Thus $\pi$ is injective on $\mathbb N^2$.

This is not a pairing compiler for arbitrary primitive recursion. A
numerical decoder is not supplied by the displayed term, and the
classification prevents the usual coordinate projections of this pairing
from being positive-iteration unary terms. For example, either projection
is nonmonotone on suitable values in the range of $\pi$, while a
nonconstant unary term is strictly increasing on all of $\mathbb N$.
For the first projection use $\pi(1,0)=11<13=\pi(0,2)$;
for the second use $\pi(0,1)=9<11=\pi(1,0)$.

\subsection{Why quotient equality is a substantive escape}

There is an explicit standard-model quotient with term-defined ring
operations. This illustrates the limits of
Theorem~\ref{thm:positive-actual-unit}, not a bi-interpretation of the
whole eager theory.
In addition to $q$, put
\[
                  m(a,b):=F(a,0;u.q(u,b))=(a+1)(b+1).
\]
On $\mathbb N^2$ define
\[
 (a,b)\mathrel E(c,d)
       \quad\Longleftrightarrow\quad q(a,d)=q(c,b).
\]
This is exactly equality of the integer differences $a-b$ and $c-d$.
The map $(a,b)\mapsto a-b$ identifies $\mathbb N^2/E$ with
$\mathbb Z$. Zero is represented by $(0,0)$, the multiplicative unit
by $(1,0)$, and negation swaps coordinates. Addition and multiplication
are given by the following
\emph{tuples of positive-iteration terms}:
\[
\begin{aligned}
 (a,b)\mathbin\oplus(c,d)
     &:=\bigl(q(a,c),q(b,d)\bigr),\\
 (a,b)\mathbin\otimes(c,d)
     &:=\bigl(q(m(a,c),m(b,d)),\ q(m(a,d),m(b,c))\bigr).
\end{aligned}
\]
For addition, the two extra ones cancel. For multiplication, expanding
the difference of the two coordinates gives
\[
 (a+1)(c+1)+(b+1)(d+1)
 -(a+1)(d+1)-(b+1)(c+1)
 =(a-b)(c-d).
\]
Thus both operations respect $E$ and induce the usual integer ring.
The nonnegative-difference domain also has the first-order description
\[
                 a=b\quad\lor\quad\exists r\,(q(b,r)=a).
\]
These are standard-algebra identities. No claim is made here that the
eager axioms prove the required quotient properties, transfer induction
to graph-defined arithmetic, or supply closure under arbitrary primitive
recursion.

In fact the inherited two-element eager model has $Sx=1$ and interprets
$F(n,a;h)$ as $h(a)$. There $q(a,b)=1$ for all inputs, so the same
formula $E$ identifies every pair and the displayed quotient has one
element. This is a concrete reason why a successful construction in
the standard algebra is insufficient for interpretation of iPRA in the
whole proposed eager theory. Adding successor injectivity excludes that
particular model, but proving a new interpretation after such an addition
would be further work.

\paragraph{Verification boundary.}
The companion Python program checks a deterministic finite sample of the
classification algorithm, and the explicit pairing and quotient formulas.
It does not verify the structural induction proof in a proof assistant.
The proofs above establish the unbounded standard-function statements;
the tests check their transcription on the reported finite inputs.
